\documentclass[11pt]{amsart}

\usepackage[T1]{fontenc}
\usepackage{lmodern}
\usepackage{microtype}
\usepackage{mathtools,amssymb,amsthm}
\usepackage[hidelinks]{hyperref}

\hypersetup{
  pdftitle={Counterexamples to the Iterated Beatty Range Criterion Below 5/4, and a Sufficient Per-Level Cutoff}
}

\newtheorem{theorem}{Theorem}
\newtheorem{lemma}[theorem]{Lemma}
\newtheorem{proposition}[theorem]{Proposition}
\newtheorem{corollary}[theorem]{Corollary}
\theoremstyle{remark}
\newtheorem*{remark}{Remark}
\newtheorem*{question}{Question 1 of \cite{KVZ}}

\DeclarePairedDelimiter{\fl}{\lfloor}{\rfloor}
\newcommand{\fp}[1]{\left\{#1\right\}}
\newcommand{\RC}{\mathrm{RC}}

\title[The iterated Beatty range criterion]
{Counterexamples to the Iterated Beatty Range Criterion Below $5/4$,
and a Sufficient Per-Level Cutoff}

\date{}

\subjclass[2020]{11B83, 11A67, 11U05}
\keywords{Beatty sequence, iterated floor function, range criterion, golden ratio,
$n$-bonacci constant, open question}

\begin{document}

\begin{abstract}
For irrational $\alpha>1$ let $f_\alpha(y)=\fl{\alpha y}$ on the positive integers.
Khani, Valizadeh and Zarei prove that for irrational $\alpha>(3+\sqrt5)/2\approx2.618$
a positive integer $x$ lies in the range of $f_\alpha^{\,n}$ if and only if an explicit
system $\RC_{\le n}(\alpha,x)$ of inequalities on
$\fp{x/\alpha},\dots,\fp{x/\alpha^{n}}$ holds, and they ask (Question~1) whether the
same equivalence holds for irrational $\alpha\in(1,2.618\ldots)$. It fails on $(1,5/4]$:
if $1<\alpha<\varphi=(1+\sqrt5)/2$ then every integer $x$ with
$\alpha^{2}/(\alpha^{2}-1)<x<1/(\alpha-1)$ is a fixed point of $f_\alpha$, hence lies in
the range of every iterate, yet $\RC_{\le 2}(\alpha,x)$ fails, and for irrational
$\alpha\le(1+\sqrt{17})/4=1.2807\ldots$, in particular for every irrational
$\alpha\in(1,5/4]$, the integer $x^{*}=\fl{1/(\alpha-1)}$ lies in that window. It does
hold at level $n$ under the sufficient condition
$\alpha^{n}\ge\alpha^{n-1}+\dots+\alpha+1$, equivalently $\alpha\ge a_n$ (the
$n$-bonacci constant); since $a_n<2$ for all $n$, it holds at every level for every
irrational $\alpha>2$. One direction, $\RC_{\le n}\Rightarrow$ membership, needs no
hypothesis on $\alpha$ at all. Sharpness of $a_n$ is not claimed: the cutoff is
sufficient only, and no per-level threshold is established here.
\end{abstract}

\maketitle

\section{Question 1 and the answer}

Throughout, $\alpha>1$ is irrational, $\fl{t}$ and $\fp{t}=t-\fl{t}$ are the integral
and fractional parts of $t$, and
\[
 f=f_\alpha:\mathbb Z_{\ge1}\to\mathbb Z_{\ge1},\qquad f(y)=\fl{\alpha y}.
\]
We write $\mathcal R_n=f^{\,n}(\mathbb Z_{\ge1})$ for the range of the $n$-th iterate.
``Natural number'' means \emph{positive} integer; this is forced by the statement below,
whose first clause reads $1-1/\alpha<\fp{x/\alpha}<1$ and so is false at $x=0$.

Khani, Valizadeh and Zarei \cite{KVZ} prove the following (their unnumbered main
theorem).

\begin{quote}
\emph{Suppose that $n\ge1$, and $\alpha$ is an irrational number strictly greater than
$\frac{3+\sqrt5}{2}$. Then, a natural number $x$ belongs to the range of $f^{n}$ if and
only if $1-\frac1\alpha<\fp{\frac x\alpha}<1$ and for each $i=2,\dots,n$ we have that}
\begin{equation}\label{eq:CRn}
 1-\frac2\alpha<\fp{\frac{x}{\alpha^{i}}}-\frac{\fp{\frac{x}{\alpha^{i-1}}}}{\alpha}
 <1-\frac1\alpha .
\end{equation}
\end{quote}

\noindent
Write $\RC_{\le n}(\alpha,x)$ for the conjunction of the first clause with
\eqref{eq:CRn} for $i=2,\dots,n$, and $\RC_i(\alpha,x)$ for its $i$-th clause. The paper
\cite{KVZ} closes with two questions; the first is reproduced here in full.

\begin{question}
\emph{Our assumption on $\alpha$ being strictly greater than
$\frac{3+\sqrt5}{2}\approx2.618$ implies that it satisfies the algebraic inequalities
below:
\[
 \alpha^{2}-3\alpha+1>0\quad\text{ and }\quad\alpha^{4}-2\alpha^{3}-\alpha+1>0.
\]
In fact, the greatest root of the polynomials above is approximately $2.618$ and
$2.118$, respectively. As we observed, these inequalities are crucial to prove Lemma
\textup{[}Technical Lemma\textup{]}. A natural question is whether our theorem holds also
for irrationals belonging to the following set:
\[
 (1,\,2.618\!\sim)=(1,\,2.118\!\sim)\cup(2.118\!\sim,\,2.618\!\sim).
\]
In this regard, recall that the function $f_\alpha$ is surjective whenever $\alpha$ is
less than or equal to one.}
\end{question}

\noindent
A single counterexample answers Question~1 in the negative; the ones produced here are
elementary. What is proved below says where the criterion fails, and gives a sufficient
condition for it to hold.

\begin{itemize}
\item \textbf{No, on $(1,5/4]$.} Theorem~\ref{thm:B} and Corollary~\ref{cor:B} produce,
 for every irrational $\alpha\in(1,5/4]$, an integer $x$ that lies in $\mathcal R_n$ for
 \emph{every} $n$ and for which $\RC_{\le n}(\alpha,x)$ is false already at index $2$.
 The witnesses are fixed points of $f_\alpha$.
\item \textbf{Yes, on $\bigl(2,(3+\sqrt5)/2\bigr)$.} Theorem~\ref{thm:A} proves the
 equivalence for every irrational $\alpha>2$ and every $n$, and at level $n$ under the
 sufficient condition $\alpha^{n}\ge\alpha^{n-1}+\dots+1$, which is not claimed to be
 necessary.
\item \textbf{Open on $(\varphi,2)$.} For $\alpha$ in that band Theorem~\ref{thm:A}
 settles the levels $n$ with $\alpha\ge a_n$ and leaves the others undecided, and we
 exhibit no counterexample there. At the left endpoint $\alpha=\varphi$ the equivalence
 itself is settled in the negative, at level $5$ (Section~\ref{sec:phi}); level $n=3$ at
 $\varphi$ is left open. See Section~\ref{sec:open}.
\end{itemize}

Nothing here contradicts the theorem of \cite{KVZ}: its hypothesis
$\alpha>(3+\sqrt5)/2$ is untouched, and both witnesses below lie at or below $\varphi<2$.

The case $n=1$ of the criterion is published as Fraenkel--Levitt--Shimshoni \cite{FLS}
and, in the short form used here, as Sz\"usz \cite{Szusz}.

\section{Two elementary lemmas}\label{sec:lemmas}

Following \cite{KVZ} put $v_0(x)=x$ and $v_k(x)=\fl{v_{k-1}(x)/\alpha}+1$, and
\[
 \rho_0(x)=0,\qquad
 \rho_k(x)=\frac{1-\fp{x/\alpha^{k}}}{\alpha}\quad(k\ge1),
\]
so that $\rho_k(x)\in(0,1/\alpha)$ for $k\ge1$, because $x/\alpha^{k}$ is irrational.
Multiplying by $\alpha$ one checks at once that for $k\ge1$
\begin{equation}\label{eq:rcform}
 \RC_k(\alpha,x)\iff
 0<\rho_k(x)-\frac{\rho_{k-1}(x)}{\alpha}<\frac1{\alpha^{2}},
\end{equation}
the case $k=1$ reading $0<\rho_1(x)<1/\alpha^{2}$: indeed
$\alpha^{2}\bigl(\rho_k-\rho_{k-1}/\alpha\bigr)
=\alpha-1-\alpha\fp{x/\alpha^{k}}+\fp{x/\alpha^{k-1}}$, and dividing the resulting
double inequality $0<\cdot<1$ by $\alpha$ gives exactly \eqref{eq:CRn}. We use the form
\eqref{eq:rcform} from now on.

\begin{lemma}\label{lem:R1}
Let $y\ge1$ be an integer. Then $y\in\mathcal R_1$ if and only if
$\rho_1(y)<1/\alpha^{2}$, and in that case the unique $z$ with $f(z)=y$ is $v_1(y)$.
\end{lemma}

\begin{proof}
$f(z)=y$ means $y\le\alpha z<y+1$, i.e.\ $y/\alpha\le z<(y+1)/\alpha$. As $\alpha$ is
irrational and $y\ge1$, $y/\alpha\notin\mathbb Z$, so an integer $z$ in that window
exists if and only if $\fl{y/\alpha}+1<(y+1)/\alpha$, and then it is unique and equals
$v_1(y)$. Finally $\fl{y/\alpha}+1<(y+1)/\alpha$ is $1-\fp{y/\alpha}<1/\alpha$, that is
$\rho_1(y)<1/\alpha^{2}$.
\end{proof}

\begin{lemma}\label{lem:tower}
For all $n\ge1$ and all integers $x\ge1$:
$x\in\mathcal R_n$ if and only if $v_k(x)\in\mathcal R_1$ for $k=0,1,\dots,n-1$.
\end{lemma}

\begin{proof}
By Lemma~\ref{lem:R1}, $x\in\mathcal R_{n}$ iff $x\in\mathcal R_1$ and
$v_1(x)\in\mathcal R_{n-1}$; iterate, using $v_k=v_1\circ v_{k-1}$.
\end{proof}

The next lemma is the whole mechanism. Say that $x$ \emph{carries at level $k$} if
$c_k(x):=\fl{\fp{x/\alpha^{k}}+\rho_{k-1}(x)}$ equals $1$; since
$\fp{x/\alpha^{k}}<1$ and $\rho_{k-1}(x)<1/\alpha<1$, indeed $c_k(x)\in\{0,1\}$.

\begin{lemma}\label{lem:key}
Fix $x\ge1$ and $k\ge1$, and suppose
$v_{k-1}(x)=x/\alpha^{k-1}+\alpha\rho_{k-1}(x)$. Then
\begin{align}
 \rho_1\bigl(v_{k-1}(x)\bigr)
 &=\rho_k(x)-\frac{\rho_{k-1}(x)}{\alpha}+\frac{c_k(x)}{\alpha},
 \label{eq:key1}\\
 v_k(x)&=\fl{x/\alpha^{k}}+1+c_k(x),
 \label{eq:key2}
\end{align}
and $c_k(x)=0$ if and only if $\rho_k(x)-\rho_{k-1}(x)/\alpha>0$. Consequently, if
$c_k(x)=1$ then $\RC_k(\alpha,x)$ is false, while if $c_k(x)=0$ then
\[
 \RC_k(\alpha,x)\iff v_{k-1}(x)\in\mathcal R_1
 \qquad\text{and}\qquad
 v_k(x)=x/\alpha^{k}+\alpha\rho_k(x).
\]
\end{lemma}

\begin{proof}
Dividing the hypothesis by $\alpha$ gives
$v_{k-1}(x)/\alpha=x/\alpha^{k}+\rho_{k-1}(x)$, whence
$\fp{v_{k-1}(x)/\alpha}=\fp{x/\alpha^{k}}+\rho_{k-1}(x)-c_k(x)$ and
\[
 \rho_1\bigl(v_{k-1}(x)\bigr)
 =\frac{1-\fp{x/\alpha^{k}}-\rho_{k-1}(x)+c_k(x)}{\alpha}
 =\rho_k(x)-\frac{\rho_{k-1}(x)}{\alpha}+\frac{c_k(x)}{\alpha},
\]
because $\alpha\bigl(\rho_k(x)-\rho_{k-1}(x)/\alpha\bigr)
=1-\fp{x/\alpha^{k}}-\rho_{k-1}(x)$. The same identity shows that $c_k(x)=0$, i.e.\
$\fp{x/\alpha^{k}}+\rho_{k-1}(x)<1$, is equivalent to
$\rho_k(x)-\rho_{k-1}(x)/\alpha>0$. Next,
\[
 v_k(x)=\fl{\tfrac{v_{k-1}(x)}\alpha}+1
 =\frac{v_{k-1}(x)}\alpha+\alpha\rho_1\bigl(v_{k-1}(x)\bigr)
 =\frac{x}{\alpha^{k}}+1-\fp{\tfrac{x}{\alpha^{k}}}+c_k(x),
\]
which is \eqref{eq:key2}; when $c_k(x)=0$ this is
$x/\alpha^{k}+\alpha\rho_k(x)$. Finally, if $c_k(x)=1$ then
$\rho_k(x)-\rho_{k-1}(x)/\alpha\le0$ and $\RC_k$ fails by \eqref{eq:rcform}, while if
$c_k(x)=0$ then $\rho_k(x)-\rho_{k-1}(x)/\alpha=\rho_1(v_{k-1}(x))$, which is positive,
so $\RC_k(\alpha,x)$ reduces to $\rho_1(v_{k-1}(x))<1/\alpha^{2}$; apply
Lemma~\ref{lem:R1}.
\end{proof}

\section{The positive half}\label{sec:A}

One implication of \eqref{eq:CRn} costs nothing.

\begin{proposition}\label{prop:free}
For every irrational $\alpha>1$, every $n\ge1$ and every integer $x\ge1$:
if $\RC_{\le n}(\alpha,x)$ holds then $x\in\mathcal R_n$.
\end{proposition}

\begin{proof}
Induct on $k=1,\dots,n$ the statement ``$v_{k-1}(x)\in\mathcal R_1$ and
$v_k(x)=x/\alpha^{k}+\alpha\rho_k(x)$''. The hypothesis of Lemma~\ref{lem:key} holds for
$k=1$ because $v_0(x)=x=x/\alpha^{0}+\alpha\rho_0(x)$, and is carried forward by the
induction. At each step $\RC_k(\alpha,x)$ gives
$\rho_k(x)-\rho_{k-1}(x)/\alpha>0$, hence $c_k(x)=0$ by Lemma~\ref{lem:key}, hence
$v_{k-1}(x)\in\mathcal R_1$ and $v_k(x)=x/\alpha^{k}+\alpha\rho_k(x)$. Now apply
Lemma~\ref{lem:tower}.
\end{proof}

So every failure of the equivalence is of one shape: $x$ in the range, criterion false.
The converse implication is where a hypothesis on $\alpha$ is needed; the one used below
is sufficient, and is not claimed to be necessary. For $n\ge1$ let
\[
 p_n(t)=t^{n+1}-2t^{n}+1 .
\]

\begin{lemma}\label{lem:an}
For each $n\ge2$, $p_n$ has exactly one root $a_n$ in $(1,\infty)$; it satisfies
$1<a_n<2$, and for $t>1$
\[
 t^{n}\ge t^{n-1}+t^{n-2}+\dots+t+1
 \iff p_n(t)\ge0
 \iff t\ge a_n .
\]
Moreover $a_2<a_3<a_4<\cdots<2$ and $a_n\to2$, with
\[
 a_2=\varphi=1.6180339887\ldots,\qquad a_3=1.8392867552\ldots,
\]
\[
 a_4=1.9275619754\ldots,\qquad a_5=1.9659482366\ldots,
\]
\[
 a_6=1.9835828434\ldots
\]
For $n=1$ the displayed inequality reads $t\ge1$ and holds for all $t>1$; we set
$a_1=1$.
\end{lemma}

\begin{proof}
Multiplying $t^{n}\ge\sum_{i=0}^{n-1}t^{i}=(t^{n}-1)/(t-1)$ by $t-1>0$ gives
$t^{n+1}-t^{n}\ge t^{n}-1$, i.e.\ $p_n(t)\ge0$; the steps are reversible. Now
$p_n'(t)=t^{n-1}\bigl((n+1)t-2n\bigr)$, so on $(1,\infty)$ the function $p_n$ decreases
on $\bigl(1,\tfrac{2n}{n+1}\bigr)$ and increases on $\bigl(\tfrac{2n}{n+1},\infty\bigr)$.
Since $p_n(1)=0$ we get $p_n<0$ on $\bigl(1,\tfrac{2n}{n+1}\bigr]$, and since
$p_n(2)=2^{n+1}-2^{n+1}+1=1>0$ there is exactly one root $a_n\in\bigl(\tfrac{2n}{n+1},2\bigr)$,
with $p_n\ge0$ on $[a_n,\infty)$ and $p_n<0$ on $(1,a_n)$. For monotonicity,
$p_{n+1}(t)-p_n(t)=t^{n}(t-1)(t-2)<0$ for $t\in(1,2)$, so $p_{n+1}(a_n)<p_n(a_n)=0$ and
therefore $a_{n+1}>a_n$. For the limit, $p_n(2-\varepsilon)=1-\varepsilon(2-\varepsilon)^{n}\to-\infty$
for fixed $\varepsilon\in(0,1)$, so $a_n>2-\varepsilon$ eventually. Finally
$p_2(t)=t^{3}-2t^{2}+1=(t-1)(t^{2}-t-1)$, so $a_2=\varphi$; the four decimals are
recomputed in the accompanying program.
\end{proof}

\begin{theorem}\label{thm:A}
Let $\alpha>1$ be irrational, let $n\ge1$, and assume
\begin{equation}\label{eq:thresh}
 \alpha^{n}\ \ge\ \alpha^{n-1}+\alpha^{n-2}+\dots+\alpha+1,
 \qquad\text{equivalently } \alpha\ge a_n .
\end{equation}
Then for every integer $x\ge1$,
\[
 x\in\mathcal R_n\iff\RC_{\le n}(\alpha,x).
\]
In particular the equivalence holds for every $n\ge1$ whenever $\alpha>2$, and at level
$n$ whenever $\alpha\ge a_n$.
\end{theorem}

\begin{proof}
Proposition~\ref{prop:free} gives $\Leftarrow$ with no hypothesis. For $\Rightarrow$
assume $x\in\mathcal R_n$; by Lemma~\ref{lem:tower}, $v_{k-1}(x)\in\mathcal R_1$ for
$k=1,\dots,n$, i.e.\ $\rho_1\bigl(v_{k-1}(x)\bigr)<1/\alpha^{2}$ by Lemma~\ref{lem:R1}.

Dividing \eqref{eq:thresh} by $\alpha^{n}$ turns it into
$1\ge\sum_{j=1}^{n}\alpha^{-j}$, i.e.
\begin{equation}\label{eq:thresh2}
 \sum_{j=2}^{k}\alpha^{-j}\ \le\ \sum_{j=2}^{n}\alpha^{-j}\ \le\ 1-\frac1\alpha
 \qquad (1\le k\le n),
\end{equation}
the first inequality because all terms are positive.

We prove by induction on $k=1,\dots,n$ that $c_k(x)=0$, that $\RC_k(\alpha,x)$ holds and
that $v_k(x)=x/\alpha^{k}+\alpha\rho_k(x)$; the last clause supplies the hypothesis of
Lemma~\ref{lem:key} at the next step, and holds for $k=0$ trivially. For $k=1$,
$c_1(x)=\fl{\fp{x/\alpha}}=0$, so Lemma~\ref{lem:key} gives everything.

Let $k\ge2$ and suppose the claim holds below $k$. Then $\RC_j(\alpha,x)$ holds for
$j\le k-1$, so by \eqref{eq:rcform} $\rho_j(x)<\alpha^{-2}+\rho_{j-1}(x)/\alpha$ for
those $j$; starting from $\rho_0(x)=0$ this telescopes to
\[
 \rho_{k-1}(x)\ <\ \sum_{j=2}^{k}\alpha^{-j}\ \le\ 1-\frac1\alpha,
\]
the last step by \eqref{eq:thresh2}. Suppose $c_k(x)=1$. Then, by \eqref{eq:key1} and
$\rho_k(x)>0$,
\[
 \rho_1\bigl(v_{k-1}(x)\bigr)
 =\rho_k(x)-\frac{\rho_{k-1}(x)}{\alpha}+\frac1\alpha
 >\frac{1-\rho_{k-1}(x)}{\alpha}
 >\frac{1-\bigl(1-\tfrac1\alpha\bigr)}{\alpha}
 =\frac1{\alpha^{2}},
\]
contradicting $v_{k-1}(x)\in\mathcal R_1$. Hence $c_k(x)=0$, so Lemma~\ref{lem:key}
turns $v_{k-1}(x)\in\mathcal R_1$ into $\RC_k(\alpha,x)$, and it also gives the identity
$v_k(x)=x/\alpha^{k}+\alpha\rho_k(x)$.

If $\alpha>2$ then $\sum_{j=1}^{n}\alpha^{-j}<\sum_{j\ge1}\alpha^{-j}=1/(\alpha-1)<1$,
so \eqref{eq:thresh} holds for every $n$; alternatively $a_n<2<\alpha$ by
Lemma~\ref{lem:an}.
\end{proof}

\section{The negative half}\label{sec:B}

\begin{theorem}\label{thm:B}
Let $\alpha$ be irrational with $1<\alpha<\varphi$, and let $x$ be an integer with
\begin{equation}\label{eq:window}
 \frac{\alpha^{2}}{\alpha^{2}-1}\ <\ x\ <\ \frac1{\alpha-1}.
\end{equation}
Then
\begin{enumerate}
\item[(i)] $f_\alpha(x)=x$; consequently $x\in\mathcal R_n$ for every $n\ge1$;
\item[(ii)] $\fl{x/\alpha}=x-1$ and $\fl{x/\alpha^{2}}=x-2$, and the quantity appearing
 in \eqref{eq:CRn} at $i=2$ is
 \[
  D_2(\alpha,x):=\fp{\frac{x}{\alpha^{2}}}-\frac{\fp{x/\alpha}}{\alpha}
  =\frac{x-1}{\alpha}-(x-2)
  =\Bigl(1-\frac1\alpha\Bigr)+\fp{\frac x\alpha};
 \]
\item[(iii)] $\RC_1(\alpha,x)$ holds and $\RC_2(\alpha,x)$ fails, on the upper side:
 $D_2(\alpha,x)>1-1/\alpha$. Hence $\RC_{\le n}(\alpha,x)$ is false for every $n\ge2$,
 while $x\in\mathcal R_n$.
\end{enumerate}
The interval in \eqref{eq:window} is nonempty if and only if $\alpha<\varphi$.
\end{theorem}

\begin{proof}
Nonemptiness: $\alpha^{2}/(\alpha^{2}-1)<1/(\alpha-1)$ iff $\alpha^{2}(\alpha-1)<\alpha^{2}-1$
iff $\alpha^{2}<\alpha+1$ iff $\alpha<\varphi$.

(i) $x<1/(\alpha-1)$ gives $\alpha x<x+1$, and $\alpha x>x$, so $\fl{\alpha x}=x$. A
fixed point of $f_\alpha$ lies in $\mathcal R_n$ for all $n$.

(ii) From $x<1/(\alpha-1)<\alpha/(\alpha-1)$ we get $x(\alpha-1)<\alpha$, i.e.
$x/\alpha>x-1$; also $x/\alpha<x$; so $\fl{x/\alpha}=x-1$. The left inequality of
\eqref{eq:window} gives $x(\alpha^{2}-1)>\alpha^{2}$, i.e.\ $x/\alpha^{2}<x-1$; and
$\alpha+1\le2\alpha^{2}$ for $\alpha\ge1$ gives
$1/(\alpha-1)=(\alpha+1)/(\alpha^{2}-1)\le2\alpha^{2}/(\alpha^{2}-1)$, so
$x(\alpha^{2}-1)<2\alpha^{2}$, i.e.\ $x/\alpha^{2}>x-2$; so $\fl{x/\alpha^{2}}=x-2$.
Therefore
\[
 D_2(\alpha,x)
 =\Bigl(\frac{x}{\alpha^{2}}-x+2\Bigr)-\frac1\alpha\Bigl(\frac x\alpha-x+1\Bigr)
 =\frac{x-1}\alpha-(x-2),
\]
and $\bigl(1-\tfrac1\alpha\bigr)+\fp{x/\alpha}
=1-\tfrac1\alpha+\tfrac x\alpha-x+1=\tfrac{x-1}\alpha-(x-2)$, the same number.

(iii) $\RC_1$ says $\fp{x/\alpha}>1-1/\alpha$, i.e.\ $x/\alpha-(x-1)>1-1/\alpha$, i.e.\
$(x+1)/\alpha>x$, i.e.\ $x(\alpha-1)<1$: true by \eqref{eq:window}. And by (ii)
\[
 D_2(\alpha,x)-\Bigl(1-\frac1\alpha\Bigr)=\fp{\frac x\alpha}>0,
\]
strictly, since $x/\alpha$ is irrational. So the right-hand inequality of
\eqref{eq:CRn} fails at $i=2$.
\end{proof}

\begin{corollary}\label{cor:B}
Let $\alpha$ be irrational with $1<\alpha\le\frac{1+\sqrt{17}}4=1.2807764064\ldots$ --
in particular let $\alpha\in(1,5/4]$ be irrational. Then
$x^{*}=\fl{1/(\alpha-1)}$ satisfies \eqref{eq:window}, so
Theorem~\ref{thm:B} applies to it: $x^{*}$ is a fixed point of $f_\alpha$ and
$\RC_{\le n}(\alpha,x^{*})$ fails for every $n\ge2$.
\end{corollary}

\begin{proof}
Put $L=1/(\alpha-1)$, which is irrational, so $x^{*}=\fl L<L$: the right half of
\eqref{eq:window} holds. Also
$\alpha^{2}/(\alpha^{2}-1)=cL$ with $c=\alpha^{2}/(\alpha+1)$, and $c<1$ because
$\alpha<\varphi$. Now $x^{*}>cL$ is $(1-c)L>\fp L$, and $\fp L<1$, so it suffices that
$(1-c)L\ge1$. But
\[
 (1-c)L=\frac{\alpha+1-\alpha^{2}}{\alpha^{2}-1},
\]
and $\alpha+1-\alpha^{2}\ge\alpha^{2}-1$ is $2\alpha^{2}-\alpha-2\le0$, i.e.\
$\alpha\le\frac{1+\sqrt{17}}4$. Finally $5/4<\frac{1+\sqrt{17}}4$.
\end{proof}

\begin{remark}
Corollary~\ref{cor:B} is a sufficient condition, not a characterisation: for
$\alpha\in\bigl(\frac{1+\sqrt{17}}4,\varphi\bigr)$ the interval \eqref{eq:window} is
nonempty but need not contain an integer, and then Theorem~\ref{thm:B} produces nothing.
\end{remark}

\subsection*{A worked witness}
Take
\[
 \alpha=\frac{7+\sqrt2}8=1.0517766952966\ldots,
 \qquad 64\alpha^{2}-112\alpha+47=0,
\]
the unique root of $64t^{2}-112t+47$
in $(1.05,1.06)$; it is irrational and less than $5/4$. Here $\alpha-1=(\sqrt2-1)/8$ and
$\alpha^{2}-1=(14\sqrt2-13)/64$, so \eqref{eq:window} is
\[
 \frac{1055+896\sqrt2}{223}\ <\ x\ <\ 8+8\sqrt2,
\]
that is,
\[
 10.4131630129\ldots\ <\ x\ <\ 19.3137084989\ldots,
\]
whose integers are $x=11,12,\dots,19$, the largest being
$x^{*}=\fl{1/(\alpha-1)}=19$. For each of them $f_\alpha(x)=x$, and by
Theorem~\ref{thm:B}(ii)
\[
 D_2(\alpha,x)=\frac{9x+38-8(x-1)\sqrt2}{47},
 \qquad
 1-\frac1\alpha=\frac{8\sqrt2-9}{47},
\]
the latter being $0.0492278404\ldots$;
so $D_2(\alpha,x)>1-1/\alpha$ is equivalent to $x<9+8\sqrt2=20.3137084989\ldots$, which
the nine values satisfy. An exhaustive check at this $\alpha$ of the equivalence
$x\in\mathcal R_2\Leftrightarrow\RC_{\le2}(\alpha,x)$ over $x=1,\dots,24$ finds it
\emph{false} at exactly these nine values of $x$ and \emph{true} at the other fifteen;
nothing here rules out further counterexamples at this $\alpha$ above $x=24$.

\section{A second witness: \texorpdfstring{$\alpha=\varphi$, $n=5$, $x=22$}{alpha=phi, n=5, x=22}}
\label{sec:phi}

The witnesses of Theorem~\ref{thm:B} all fail at level $2$ and all have
$\alpha<\varphi$. The equivalence also fails higher up, at $\alpha$ equal to the golden
ratio itself, where Theorem~\ref{thm:A} guarantees levels $1$ and $2$ (because
$\varphi=a_2$) and permits failure from level $3$ on.

Let $\alpha=\varphi=(1+\sqrt5)/2$, so $1/\varphi=\varphi-1$. Then
\[
 3\xrightarrow{\ f\ }4\xrightarrow{\ f\ }6\xrightarrow{\ f\ }9
 \xrightarrow{\ f\ }14\xrightarrow{\ f\ }22,
\]
each step being $y\mapsto\fl{\varphi y}$, so $22\in\mathcal R_5$. On the other hand, from
$\varphi^{-1}=\varphi-1$, $\varphi^{-2}=2-\varphi$, $\varphi^{-3}=2\varphi-3$,
$\varphi^{-4}=5-3\varphi$, $\varphi^{-5}=5\varphi-8$ one computes in $\mathbb Z[\varphi]$
\[
 \fp{22/\varphi}=22\varphi-35,\quad
 \fp{22/\varphi^{2}}=36-22\varphi,\quad
 \fp{22/\varphi^{3}}=44\varphi-71,
\]
\[
 \fp{22/\varphi^{4}}=107-66\varphi,\quad
 \fp{22/\varphi^{5}}=110\varphi-177,
\]
and hence, writing $D_i=\fp{22/\varphi^{i}}-\fp{22/\varphi^{i-1}}/\varphi$,
\begin{align*}
 D_2&=13\varphi-21=\phantom{-}0.0344418537\ldots, &
 D_3&=8\varphi-13=-0.0557280900\ldots,\\
 D_4&=5\varphi-8=\phantom{-}0.0901699437\ldots, &
 D_5&=3\varphi-4=\phantom{-}0.8541019662\ldots
\end{align*}
The bounds of \eqref{eq:CRn} are $1-2/\varphi=3-2\varphi=-0.2360679774\ldots$ and
$1-1/\varphi=2-\varphi=0.3819660112\ldots$, so $D_2,D_3,D_4$ satisfy \eqref{eq:CRn}
while
\[
 D_5-\Bigl(1-\frac1\varphi\Bigr)=(3\varphi-4)-(2-\varphi)=4\varphi-6
 =0.4721359549\ldots>0 .
\]
Thus $\RC_{\le4}(\varphi,22)$ holds, $\RC_{\le5}(\varphi,22)$ fails at $i=5$, and
$22\in\mathcal R_5$: the equivalence fails at $\alpha=\varphi$, $n=5$, $x=22$.

\begin{remark}
The identity that supplies a closed form for $f_\varphi^{\,n}$,
namely $f_\varphi^{\,2}=f_\varphi+\mathrm{id}-1$, is in print: Theorem~1 of Connell
\cite{Connell}, a lemma of Khani and Zarei \cite{KhaniZarei}, and axiom (A1.4)
of \cite{KVZ-Fib}.
\end{remark}

\section{What is \emph{not} settled}\label{sec:open}

For $\alpha$ in the band $(\varphi,2)$ we exhibit no counterexample, and
Theorem~\ref{thm:A} settles level $n$ only when $\alpha\ge a_n$, leaving the levels $n$
with $a_n>\alpha$ undecided in both directions. Since $\varphi=a_2$,
Theorem~\ref{thm:A} covers levels $1$ and $2$ at $\alpha=\varphi$ and permits failure
from level $3$ on; Section~\ref{sec:phi} exhibits a failure at level $5$, and whether
level $3$ fails at $\varphi$ we do not know. We do not claim that $a_n$ is the least
$\alpha$ for which the conclusion of Theorem~\ref{thm:A} holds at level $n$: condition
\eqref{eq:thresh} is sufficient only, so this note establishes no per-level threshold,
and the behaviour of the equivalence just below $a_n$ is among what is left open. At
$\alpha=(7+\sqrt2)/8$ the nine integers $11,\dots,19$ are exactly the disagreements for
$x\le24$ at $n=2$, and exactly the integers of \eqref{eq:window}; nothing here rules out
counterexamples above $x=24$ that are not fixed points.

\section{Verification}

Every numerical assertion above is stated exactly, as an element of
$\mathbb Q(\sqrt2)$, $\mathbb Q(\sqrt5)$ or $\mathbb Q$, and each is re-derived by
the program \texttt{verify.py} accompanying this note: each modulus is given by its
minimal polynomial and an isolating interval, and every inequality between algebraic
numbers is decided by integer arithmetic, with no floating-point comparison anywhere.
The program re-checks the equivalence of \eqref{eq:CRn} with the $\rho$-form
\eqref{eq:rcform}; both witnesses in full, including membership of $22$ in
$\mathcal R_5$ by a preimage sieve independent of the printed chain, the closed forms,
the window \eqref{eq:window} with its nine integers, and the exhaustive table $x\le24$;
the constants $a_2,\dots,a_6$ and the equivalent forms of \eqref{eq:thresh};
Theorem~\ref{thm:B} over a family of quadratic irrationals of $(1,\varphi)$; and
Theorem~\ref{thm:A} with Proposition~\ref{prop:free} by exhaustive checks of both
directions over bounded ranges of $x$ and $n$ at moduli on both sides of every cutoff
$a_n$. Those checks are bounded and a bounded check is not a proof; the theorems are
carried by their proofs above. The program states in its own output what it does not
cover.

\begin{thebibliography}{9}

\bibitem{KVZ}
M.~Khani, M.~Valizadeh, and A.~Zarei,
\emph{A note on iterated Beatty sequences},
arXiv:2607.12817v1, 14 July 2026.
\href{https://arxiv.org/abs/2607.12817}{arXiv:2607.12817}.
Statements are quoted from the \texttt{v1} source package.

\bibitem{Connell}
I.~G. Connell,
\emph{Some properties of Beatty sequences I},
Canad. Math. Bull. \textbf{2} (1959), no.~3, 190--197.

\bibitem{KhaniZarei}
M.~Khani and A.~Zarei,
\emph{The additive structure of integers with the lower Wythoff sequence},
Arch. Math. Logic \textbf{61} (2022), 1--13.
\href{https://doi.org/10.1007/s00153-022-00846-2}{doi:10.1007/s00153-022-00846-2}.

\bibitem{KVZ-Fib}
M.~Khani, M.~Valizadeh, and A.~Zarei,
\emph{Fibonacci numbers and model-complete axiomatization of Presburger arithmetic
expanded with a Beatty sequence},
arXiv:2508.02303 [math.LO], 2025.
\href{https://arxiv.org/abs/2508.02303}{arXiv:2508.02303}.

\bibitem{FLS}
A.~S. Fraenkel, J.~Levitt, and M.~Shimshoni,
\emph{Characterization of the set of values $f(n)=[n\alpha]$, $n=1,2,\ldots$},
Discrete Math. \textbf{2} (1972), 335--345.

\bibitem{Szusz}
P.~Sz\"usz,
\emph{On a theorem of Fraenkel, Levitt and Shimshoni},
Discrete Math. \textbf{56} (1985), 75--77.
\href{https://doi.org/10.1016/0012-365X(85)90194-3}{doi:10.1016/0012-365X(85)90194-3}.
Cited here from its Zentralblatt review (Zbl 0583.10002), which states the criterion in
the form used above.

\end{thebibliography}

\end{document}
